% Pursuing Stacks (cotes 134-2 to 134-8): archive cote 134-5, site page 106 (PDF page 107), typescript page 293.
% typescript. Draft compared: ps5.tex:102-159.
% AI-transcribed by Claude Opus 5.5 (claude-opus-5-5, adaptive thinking, effort high), Message Batches API, batch ps2 (msgbatch_01NVwr1Yuo8qJf5rykeGBXPK), 2026-11-12; not yet checked by eye.
\scanpage{134-5/106}
\mspage{293}

a certain feeling of uneasiness with respect to asphericity structures,
which hav'nt got yet a reasonable notion of morphism. Thus I have left
unanswered the questions raised in section 84 around the inclusion \inserted{condition}
\[ f_{!}(M_{a}) \hand{\subset} M'_{a} \qquad , \]
so it is not impossible that, while relying on the mere feeling that
the inclusion is just not reasonable and that the answer to the specific
questions are presumably negative, I am about to miss some unexpected \hand{im-}
\over{\illegible{1}}{portant} fact ! It seems I developped kind of a block against checking -
maybe the answers are well known and Tim Porter will tell me ... Maybe
I better leave the question for a later moment, as it will ripen by it-%
self if I leave it alone ...

2) I should at last introduce contractors, and morphisms of such.
When I set out on part IV of the notes, I expected that the notion of a
contractor would be one main notion, alongside with the notion of a
canonical modelizer - it turns out that so far I did'nt have much use
yet for one or the other. Contractors c\hand{\unsure{-}}an be viewed as \xout{\unsure{small}} categories
generating contractibility structures, just as aspherators are there
for generating asphericity structures. H\unsure{o}wever, whereas any small cate-%
gory is an aspherator, the same is definitely not so for contractors,
as we demand that every object in C should be contractible, for the
homotopy structure in C\^{} generated by C itself. If we except the case
of a contractor equivalent to the final category (\xout{the} a so-called \uline{trivial}
contractor), a contractor is a strict test category - thus the notion
appears somewhat as a hinge notion between the test notions, and the
\xout{x}``pure'' homotopy notions and more specifically, contractibility structu-%
res. Writing up some scratchnotes I got should be a pure routine matter.

3) A lot more interesting seems to me to try and resolve a persis-%
tent feeling of uneasiness which has been floating, throughout the
long-winded reflections on homotopy and asphericity structures in parts
III and IV of the notes. This is tied up with this fact, that \unsure{my} treat-%
ment of the main notions, namely contractibility and asphericity, has
been consistently \uline{non} \uline{auto}-\uline{dual}. More specifically, when a category M
is endowed with either a contractibility or an asphericity structure,
it does \uline{not} follow that the opposit category M$^{\circ}$ is too in a natural way.
When defining homotopy relations and homotopism structures (section 51,
A) and B)), these were autodual notions, but the notion of a\xout{n} homotopy
interval structure, which we used in order to pass from a contractibility

\margin{106}
